\documentclass[11pt,a4paper]{article}
\usepackage[utf8]{inputenc}
\usepackage[T1]{fontenc}
\usepackage[margin=0.72in]{geometry}
\usepackage{enumitem}
\usepackage{titlesec}
\usepackage{xcolor}
\usepackage{hyperref}

\definecolor{primarycolor}{RGB}{33,37,41}
\definecolor{secondarycolor}{RGB}{0,102,153}
\definecolor{accentcolor}{RGB}{95,105,115}
\hypersetup{colorlinks=true,urlcolor=secondarycolor,linkcolor=secondarycolor}
\titleformat{\section}{\large\bfseries\color{primarycolor}}{}{0em}{}
\titlespacing*{\section}{0pt}{1.1em}{0.45em}
\setlist[itemize]{leftmargin=1.2em,itemsep=0.18em,topsep=0.2em}
\pagestyle{empty}

\begin{document}
\begin{center}
{\LARGE\bfseries\color{primarycolor} MARTA KOWAL-BRANDT}\\[0.25em]
{\large QA Automation Engineer | CI/CD \& Platform Engineering}\\[0.25em]
{\small marta.kowalbrandt@example.com | +48 555 303 303 | gitlab.com/mkb-example}
\end{center}

\section*{PROFESSIONAL SUMMARY}
QA automation engineer with 7 years of experience, including hands-on ownership of CI/CD pipelines, build tooling and test environments. Built GitLab CI workflows and Helm-based ephemeral environments for three product teams, and developed a personal k3s/GitOps homelab with Terraform-provisioned Hetzner infrastructure. Seeking a DevOps or platform engineering role.

\section*{TECHNICAL SKILLS}
CI/CD and containers: GitLab CI, Jenkins, Docker, Docker Compose\\
Infrastructure and orchestration: Helm, Kubernetes (shared-cluster namespace-level access; personal k3s), Terraform (personal project), Argo CD\\
Observability and systems: Prometheus, Grafana, Loki, Linux, Git\\
Programming and data: Python, Bash, Java (test code), SQL, PostgreSQL\\
Testing: pytest, Playwright, Selenium, Postman\\
Cloud learning: AWS (EC2 and S3 basics)\\
Additional: Ansible (basic), TLS certificates with cert-manager

\section*{PROFESSIONAL EXPERIENCE}
\textbf{DevOps Engineer | Odrava Health Software} \hfill \textit{2021--Present}
\begin{itemize}
\item Own GitLab CI pipelines for three product teams (approximately 25 developers), covering build, unit, integration and end-to-end test stages and artifact publishing.
\item Reduced average pipeline time from 38 to 14 minutes by parallelising test jobs, caching dependencies and moving end-to-end tests to on-demand ephemeral environments; monthly CI runner costs fell by about 30\%.
\item Built short-lived merge-request test environments using Docker Compose and Helm on the company's shared Kubernetes cluster.
\item Maintain a pytest and Playwright end-to-end suite of approximately 900 tests; reduced the flaky-test rate from 7\% to below 1\%.
\item Wrote Bash and Python tooling to seed test data into PostgreSQL.
\item Authored internal documentation on how CI works, used for onboarding; mentored two QA interns in 2023.
\end{itemize}

\textbf{QA Engineer | Brightmill Software House} \hfill \textit{2019--2021}
\begin{itemize}
\item Performed manual and automated testing of web and mobile applications for client projects.
\item Introduced a Selenium and Java automated regression suite for the largest client project.
\item Configured Jenkins jobs for nightly regression runs.
\end{itemize}

\textbf{Customer Support Specialist (2nd line) | Brightmill Software House} \hfill \textit{2018--2019}
\begin{itemize}
\item Investigated customer issues in logs and databases and reproduced bugs for developers.
\end{itemize}

\section*{SELECTED PROJECTS}
\textbf{Homelab k3s Cluster with GitOps | Personal} \hfill \textit{2024--2025}
\begin{itemize}
\item Built a three-node Raspberry Pi 5 k3s cluster, with applications deployed from a Git repository using Argo CD.
\item Set up Prometheus, Grafana and Loki for monitoring, plus Traefik ingress and cert-manager for TLS.
\item Provisioned infrastructure for a small Hetzner VPS with Terraform; used Ansible at a basic level.
\end{itemize}

\textbf{Merge-Request Ephemeral Environments | Odrava Health} \hfill \textit{2023--2024}
\begin{itemize}
\item Implemented short-lived environments for each merge request, with the application, PostgreSQL and seeded data, deployed via Helm to the shared Kubernetes cluster from GitLab CI.
\item Configured automatic cleanup after merge or 48 hours of inactivity; QA and product teams use the environments for feature reviews, with end-to-end tests running against them.
\end{itemize}

\textbf{Flaky Test Dashboard | Odrava Health} \hfill \textit{2023}
\begin{itemize}
\item Collected CI test results in PostgreSQL and created Grafana dashboards showing flaky tests by team.
\item Teams used the dashboard to identify and fix top offenders; the flaky-test rate fell from 7\% to below 1\%.
\end{itemize}

\section*{EDUCATION AND CERTIFICATIONS}
BEng Computer Engineering | Silesian Lowlands University of Technology | 2018\\
ISTQB Foundation Level | 2019\\
Certified Kubernetes Administrator (CKA): studying; exam scheduled for January 2027 (not yet passed)

\section*{LANGUAGES}
Polish (native), English (C1), German (A2)
\end{document}